\section{Conclusions}

We presented Background-Guided Residual Flow Matching (BG-RFM) for multi-constraint seismic velocity model building when processed/interpreted constraints are already available. The method learns role-aware conditions from RMS velocity, post-stack migrated images, interpreted horizons, and sparse wells, predicts a deterministic background velocity, and applies conditional Flow Matching to the correction relative to that prediction.

The central design is prediction consistency: the same predicted background is reused to define the residual target, condition residual transport, and compose the final velocity estimate. This shared reference makes the deterministic/generative responsibility allocation explicit and gives the residual path a direct full-field interpretation as transport from a background-centered stochastic base toward the target.

Under the common 33,600-record evaluation, BG-RFM achieves RMSE 0.02695 and SSIM 0.99165 with 10.63M trainable inference parameters. The supporting formulation comparison is consistent with improved reconstruction for the complete background-guided residual organization relative to simpler full-field Flow-Matching alternatives. Together, the results support a parameter-compact and structured way to integrate heterogeneous geophysical constraints for the declared processed-constraint reconstruction task.
